%% 自動生成（tools/mkchecklist.py）。各章の「ふりかえり」の項目を集めたもの
\begingroup\footnotesize
\def\bkcb{\tikz[baseline=-0.5ex]\draw[bkTheme, line width=0.7pt, rounded corners=0.5pt] (0,-0.3em) rectangle (0.8em,0.5em);}
\begin{longtblr}{colspec={Q[c,m,wd=9mm] X[l,m] Q[c,m,wd=10mm] Q[c,m,wd=9mm] Q[c,m,wd=9mm]}, width=\linewidth, rowhead=1, colsep=3pt,
  row{1}={bg=bkVol!14, font=\sffamily\bfseries\scriptsize}, hlines={0.4pt,bkLine}, rowsep=0.7pt}
章 & できるようになったこと & ページ & 1回め & 2回め \\
{\sffamily\bfseries\color{bkTheme}5-1} & 平均の速さを $\tfrac{\Delta x}{\Delta t}$ で求められる。 & \pageref{bk:5-1:heikin} & \bkcb & \bkcb \\
 & 平均の速さが、2点を結ぶ直線の傾きだと説明できる。 & \pageref{bk:5-1:katamuki} & \bkcb & \bkcb \\
 & あいだを短くして、瞬間の速さを見積もれる。 & \pageref{bk:5-1:shunkan} & \bkcb & \bkcb \\
{\sffamily\bfseries\color{bkTheme}5-2} & $\Delta t$ を小さくして、平均の速さを表で計算できる。 & \pageref{bk:5-2:hyou} & \bkcb & \bkcb \\
 & 平均の速さを $h$ の式で表せる。 & \pageref{bk:5-2:h} & \bkcb & \bkcb \\
 & $x=at^2$ の瞬間の速さを求められる。 & \pageref{bk:5-2:a} & \bkcb & \bkcb \\
{\sffamily\bfseries\color{bkTheme}5-3} & $\lim$ の記号を読み、極限の値を求められる。 & \pageref{bk:5-3:lim} & \bkcb & \bkcb \\
 & 微分係数を、接線の傾きとして求められる。 & \pageref{bk:5-3:bibunkeisu} & \bkcb & \bkcb \\
 & 導関数を求め、接線の式を書ける。 & \pageref{bk:5-3:shiki} & \bkcb & \bkcb \\
{\sffamily\bfseries\color{bkTheme}5-4} & 表の結果から、微分のきまりを見つけられる。 & \pageref{bk:5-4:kimari0} & \bkcb & \bkcb \\
 & $(x^n)'=nx^{n-1}$ を使って微分できる。 & \pageref{bk:5-4:kimari} & \bkcb & \bkcb \\
 & 定数倍と和のきまりで、多項式を微分できる。 & \pageref{bk:5-4:wa} & \bkcb & \bkcb \\
{\sffamily\bfseries\color{bkTheme}5-5} & 速度と速さのちがいを説明できる。 & \pageref{bk:5-5:sokudo} & \bkcb & \bkcb \\
 & 位置を2回微分して、加速度を求められる。 & \pageref{bk:5-5:nikai} & \bkcb & \bkcb \\
 & $\dot{x}$、$\ddot{x}$、$\tfrac{\dd^2x}{\dd t^2}$ を読める。 & \pageref{bk:5-5:nikai} & \bkcb & \bkcb \\
{\sffamily\bfseries\color{bkTheme}5-6} & 積の微分を、長方形の面積で説明できる。 & \pageref{bk:5-6:seki} & \bkcb & \bkcb \\
 & 合成関数の内側と外側を見分けられる。 & \pageref{bk:5-6:gousei} & \bkcb & \bkcb \\
 & 合成関数の微分を使って計算できる。 & \pageref{bk:5-6:chain} & \bkcb & \bkcb \\
{\sffamily\bfseries\color{bkTheme}5-7} & 導関数の符号から、上がる・下がるを判断できる。 & \pageref{bk:5-7:fugou} & \bkcb & \bkcb \\
 & 増減表をかいて、極大・極小を求められる。 & \pageref{bk:5-7:zougen} & \bkcb & \bkcb \\
 & 範囲のはしもくらべて、最大・最小を求められる。 & \pageref{bk:5-7:saidai} & \bkcb & \bkcb \\
{\sffamily\bfseries\color{bkTheme}5-8} & 次の位置を、今の位置＋速度×$\Delta t$ で予想できる。 & \pageref{bk:5-8:tsugi} & \bkcb & \bkcb \\
 & 接線を使って、1次近似ができる。 & \pageref{bk:5-8:sessen} & \bkcb & \bkcb \\
 & 測定のずれが、計算にひびく大きさを見積もれる。 & \pageref{bk:5-8:gosa} & \bkcb & \bkcb \\
{\sffamily\bfseries\color{bkTheme}5-9} & 速さが変わるときも、細かく分けて、かけて、たして距離を求められる。 & \pageref{bk:5-9:wakeru} & \bkcb & \bkcb \\
 & 長方形・三角形・台形の面積で、進んだ距離がわかる。 & \pageref{bk:5-9:sankaku} & \bkcb & \bkcb \\
 & マイナスの面積の意味（変位と道のり）を説明できる。 & \pageref{bk:5-9:minus} & \bkcb & \bkcb \\
{\sffamily\bfseries\color{bkTheme}5-10} & 細い長方形の和で、面積を近似できる。 & \pageref{bk:5-10:hasami} & \bkcb & \bkcb \\
 & $\Sigma$ の記号を読み書きできる。 & \pageref{bk:5-10:sigma} & \bkcb & \bkcb \\
 & 定積分の意味を説明し、$x^2$ の定積分を求められる。 & \pageref{bk:5-10:teisekibun} & \bkcb & \bkcb \\
{\sffamily\bfseries\color{bkTheme}5-11} & 面積のふえる速さが高さであることを説明できる。 & \pageref{bk:5-11:kihon} & \bkcb & \bkcb \\
 & 不定積分を $+C$ をつけて書ける。 & \pageref{bk:5-11:genshi} & \bkcb & \bkcb \\
 & 定積分を $F(b)-F(a)$ で計算できる。 & \pageref{bk:5-11:teisekibun} & \bkcb & \bkcb \\
{\sffamily\bfseries\color{bkTheme}5-12} & 変わる力の仕事を、積分で求められる。 & \pageref{bk:5-12:sekibun} & \bkcb & \bkcb \\
 & 流量や電力から、たまった量を求められる。 & \pageref{bk:5-12:ryuuryou} & \bkcb & \bkcb \\
 & 積分を使って、平均の値を求められる。 & \pageref{bk:5-12:heikin} & \bkcb & \bkcb \\
{\sffamily\bfseries\color{bkTheme}5-13} & 初期条件から、積分定数を決められる。 & \pageref{bk:5-13:shoki} & \bkcb & \bkcb \\
 & 加速度から、速度と位置を求められる。 & \pageref{bk:5-13:toukasoku} & \bkcb & \bkcb \\
 & 「いまの量＝はじめの量＋変化の積分」を使える。 & \pageref{bk:5-13:teisekibun} & \bkcb & \bkcb \\
{\sffamily\bfseries\color{bkTheme}5-14} & $e$ が $(1+\frac{1}{n})^n$ の近づく先であることを確かめられる。 & \pageref{bk:5-14:n} & \bkcb & \bkcb \\
 & $(e^{ax})'=ae^{ax}$ を使える。 & \pageref{bk:5-14:eax} & \bkcb & \bkcb \\
 & $\ln x$ を微分・積分できる。 & \pageref{bk:5-14:lnbibun} & \bkcb & \bkcb \\
{\sffamily\bfseries\color{bkTheme}5-15} & 円運動の速度で、$\sin$・$\cos$ の微分を説明できる。 & \pageref{bk:5-15:en} & \bkcb & \bkcb \\
 & $\sin\omega t$、$\cos\omega t$ を微分・積分できる。 & \pageref{bk:5-15:omega} & \bkcb & \bkcb \\
 & 単振動の $\ddot{x}=-\omega^2x$ を確かめられる。 & \pageref{bk:5-15:nikai} & \bkcb & \bkcb \\
 & $\sin^2$ を積分して、交流の実効値を求められる。 & \pageref{bk:5-15:jikko} & \bkcb & \bkcb \\
{\sffamily\bfseries\color{bkTheme}5-16} & 置換積分で、$\dd u$ への書きかえができる。 & \pageref{bk:5-16:chikan} & \bkcb & \bkcb \\
 & 定積分の置換で、はんいを書きかえられる。 & \pageref{bk:5-16:hanni} & \bkcb & \bkcb \\
 & 部分積分で、$f$ と $g'$ を選んで計算できる。 & \pageref{bk:5-16:bubun} & \bkcb & \bkcb \\
{\sffamily\bfseries\color{bkTheme}5-17} & 2次近似で、曲がり方まで合わせる意味を説明できる。 & \pageref{bk:5-17:niji} & \bkcb & \bkcb \\
 & $e^x$・$\sin x$・$\cos x$ の展開を書き、値を計算できる。 & \pageref{bk:5-17:sin} & \bkcb & \bkcb \\
 & $\sin\theta\approx\theta$ から、振り子の周期を求められる。 & \pageref{bk:5-17:furiko} & \bkcb & \bkcb \\
{\sffamily\bfseries\color{bkTheme}5-18} & ふるまい表を使って、微分・積分できる。 & \pageref{bk:5-18:hyo} & \bkcb & \bkcb \\
 & 「自分にもどる」関数を答えられる。 & \pageref{bk:5-18:modoru} & \bkcb & \bkcb \\
 & はじめの状態から、式の中の数を決められる。 & \pageref{bk:5-18:shoki} & \bkcb & \bkcb \\
\end{longtblr}
\endgroup
